\begin{abstract}
\ifshort
Evaluating whether a tool-using agent can be induced to exfiltrate sensitive data
currently requires judging, per episode, whether an attack succeeded. We replace
that judgement with a construction: a synthetic record no legitimate task
transmits is planted in the agent's data store, and its bytes appearing in an
outbound tool argument addressed to an attacker-controlled host is exfiltration
by construction---sound because the adversary has unbounded freedom over the
injected instruction and none over the stolen record, which must arrive intact to
have value. Testing the obvious form of this label against a matched benign
control arm falsified it: ``the record appears in any outbound argument'' fired
on \llamaControlAnyEgress{} of \llamaControlN{} adversary-free episodes for
\llamaLabel{}---\llamaControlFPRate{} false positives---because transmitting
member records is plausibly the agent's job. Scoping to the attacker's host gave
\llamaControlExfil{}/\llamaControlN{} with detection unchanged. Any label of the
form ``a sensitive artifact appeared in an outbound channel'' inherits the
agent's base rate of legitimate transmission. Separately, with no adversary
present, \llamaLabel{} fabricated internal-sounding endpoints and posted the
planted identifier to them under a system prompt naming those records as
protected health information. Injected exfiltration was \llamaExfil{}/\llamaN{}
and \geminiExfil{}/\geminiN{}; we decline to read that gap as a claim about
deployment category, because the design confounds license with tier.
\else
Evaluating whether a tool-using language-model agent can be induced to exfiltrate
sensitive data currently requires deciding, per episode, whether an attack
succeeded---a judgement made by human transcript review or by a model-based
judge, neither of which is reproducible at byte granularity. We replace that
judgement with a construction. A synthetic record that no legitimate task needs
to transmit is planted in the agent's data store; the appearance of that
record's bytes in an outbound tool argument addressed to an attacker-controlled
host is exfiltration by construction. The construction holds because of an
asymmetry available to the defender: an adversary has unbounded freedom over the
wording of an injected instruction, and none over the stolen record, which must
arrive intact to have value.

Building the obvious version of this label and then testing it against a matched
benign control arm falsified it. The predicate ``the planted record appears in
any outbound argument'' fired on \llamaControlAnyEgress{} of \llamaControlN{}
adversary-free episodes for \llamaLabel{}---a \llamaControlFPRate{}
false-positive rate---because transmitting member records is plausibly the
agent's job. Scoping the predicate to the attacker's host moved that to
\llamaControlExfil{}/\llamaControlN{} with the injected rate unchanged. We
report this as the paper's primary result: any label of the form ``a sensitive
artifact appeared in an outbound channel'' inherits the agent's base rate of
legitimate transmission, and for an agent whose function is handling such
artifacts that rate is not small.

A second finding is independent of injection entirely. On benign tickets with no
adversary present, \llamaLabel{} fabricated internal-sounding endpoints and
posted the planted member identifier to them under a system prompt that named
those records as protected health information. In a pilot across \nModelsWord{}
deployed models and \pooledEpisodes{} episodes, injected exfiltration to the
attacker host was \llamaExfil{}/\llamaN{} and \geminiExfil{}/\geminiN{}; we
decline to read that gap as a statement about deployment category, because the
design confounds license with model tier. All \nLeaksTotal{} detected leaks were
cleartext.
\fi
\end{abstract}
